% Source: https://artofproblemsolving.com/community/c3h1480949_amc_10_2010_problem_18
% Stable identity: aops-topic-1480949-q01
\textbf{Case 1: Bernardo picks $9$.} This has probability $\dfrac{\binom82}{\binom93}=\dfrac{28}{84}=\dfrac{1}{3}$. His number then starts with $9$ and Silvia's cannot, so he wins.

\textbf{Case 2: Bernardo does not pick $9$} (probability $\dfrac{2}{3}$). Then both choose a $3$-element subset of $\{1,\dots,8\}$ uniformly at random, independently. Each has $\binom83=56$ choices, and different subsets give different numbers. So the two numbers are equal with probability $\dfrac{1}{56}$.

By symmetry (the two players are now in identical situations), ``Bernardo's number is larger'' and ``Silvia's number is larger'' have the same probability $q$. The three possibilities (larger, smaller, equal) are exclusive and cover everything, so
\[
q+q+\dfrac{1}{56}=1,\qquad q=\dfrac{1-\dfrac{1}{56}}{2}=\dfrac{55}{112}.
\]
This is why the formula has this form. The value $\dfrac{1}{2}$ is not the probability of ``Bernardo larger'' before ties are removed, so it is wrong to subtract the tie probability from it. Instead, the tie probability $\dfrac{1}{56}$ is removed from the total $1$, and the rest is shared equally. ($\dfrac{1}{2}-\dfrac{1}{56}$ would subtract the whole tie probability from one side only.)

\textbf{Total.}
\[
\dfrac{1}{3}+\dfrac{2}{3}\cdot\dfrac{55}{112}=\dfrac{56}{168}+\dfrac{55}{168}=\dfrac{111}{168}=\boxed{\dfrac{37}{56}}.
\]
